\section{Data, treatment and outcomes}

In order to conduct my analysis, I use data from \glsxtrfull{ENOE} conducted by \glsxtrfull{INEGI}. Since 2005, this survey reports quarterly data related to labour statistics for the Mexican population aged 15 and older. It is a rotating panel that follows 5 subsamples for up to five quarters with a probabilistic design.

Because of the COVID-19 pandemic, \glsxtrshort{INEGI} was not able to conduct its regular survey during the second quarter of 2020. This led to the construction of a new temporary survey (\glsxtrfull{ETOE}). The discontinuity resulted in the creation of \glsxtrfull{ENOEN}, a version of \glsxtrshort{ENOE} featuring an updated survey design and a mixture of face-to-face and telephone interviews, implemented from the third quarter of 2020 through the fourth quarter of 2022. From the first quarter of 2023 onwards, \glsxtrshort{INEGI} resumed the regular \glsxtrshort{ENOE} while maintaining comparability with \glsxtrshort{ENOEN}. To ensure consistency over time within the survey framework, I limit the sample to observations from the third quarter of 2020 through the second quarter of 2026 . Following standard \glsxtrshort{ENOE} procedures, I restrict the sample to individuals with completed interviews, valid residence status and aged 18 to 64 years old.

The \glsxtrshort{ENOE} data does not allow to follow individuals over long periods of time. Therefore, for my estimations, I aggregate the individual level data to the occupation level. The main advantage related to this data structure is that it allows me to follow cells for several quarters while also being able to measure the mass of employed people in each occupation across time. Occupations are identified using \glsxtrfull{SINCO} codes present in both \glsxtrshort{ENOEN} and \glsxtrshort{ENOE}\footnote{The \glsxtrshort{ENOE} survey uses \glsxtrshort{SINCO} codes from the 2011 methodology until the second quarter of 2021. From the third quarter of 2021 onwards the survey uses the 2019 \glsxtrshort{SINCO} classification. Because of this I harmonize codes using the official \glsxtrshort{INEGI} cross walk table.}. For each occupation $j$ in quarter $t$, the outcomes variables are defined as the weighted average outcome among workers employed in occupation $j$\footnote{Weights are specifically the \glsxtrshort{ENOE} expansion factors. This way each mean is representative of the number of workers in that occupation}. The resulting dataset is a pseudo-panel following 422 occupations\footnote{The original count is of 483 codes. However, I exclude code 9999 (non classified occupations) and restrict my analysis to occupations that have at least one worker observation across all quarters.} which represents roughly 52.6 million workers based on the last quarter of 2022. The range of weighted observations goes from 46.6 to 54 million, and raw observations go from 110 thousand to 184 thousand. This huge gap is caused by the first four quarters in my sample. Because of Covid, a large number of intervies were not ablet o be finished in person so the count drop spikes the standard errors in the first few quarters. This poses no threat to vailidity since further quarter have more stable counts and the drop is further from the treatment period.

The outcomes of interest for my analysis are the share of informal workers, the number of employed people and average wages in occupation $j$. The share of informal workers include workers employed in the informal sector as well as workers whose employment relationship lacks access to the legal and institutional protections associated with formal employment. The number of employed people come directly from the number of workers in occupation $j$ at quarter $t$.

Wages are monthly earnings in nominal pesos. However, a problem with \glsxtrshort{ENOE} is that a non random share of paid employed workers report zero hourly income despite working a positive number of hours and reporting an occupation. \cite{imput}
I treat these as missing and impute them with a predictive mean matching (PMM) method develoepdd by \cite{conasami2023imputacion}. I split workers who need imputed wages into those that did not gave an answer to a minimum wage bracket and those who did. Each group of workers is then imputed a wage using two different models. In both models, donors are workers who reported an exact amount. Then, for model 1, regresses income on worker characteristics (age group, formality, rurality, full-time status, sex, marital status, education, two-digit occupation and geographic zone) and the reported bracket. Model 2 replicates the regression without the reported minimum wage bracket. I draw the parameters and compute predicted income for all workers, and assign each worker needing imputation the observed income of one of the five closest predicted values. This is repeated ten times, and the imputed wage used in the analysis is the average across imputations.

Treatment is defined using an exposure measure to \glsxtrfull{AI} taken from \cite{martinez2024risk}. The authors use methodology from \cite{eloundou2023gptsgptsearlylook}, adapted to the Mexican context\footnote{The main assumptions in the Mexican context include: limited digital access, incomplete data, lack of infrastructure and lagged technology adoption}. Each occupation consists of a set of tasks. These tasks are subsequently categorized by human raters and GPT-4o into three categories depending on the degree to which \glsxtrshort{AI} can reduce the time required to complete them while maintaining quality. Tasks in category $E0$ are not meaningfully affected by AI capabilities, tasks in $E1$ can be heavily assisted or performed completely by \glsxtrshort{AI}, and tasks in $E2$ may become exposed to \glsxtrshort{AI} only if additional software is required to perform said task.\footnote{Details on each category defined by \cite{eloundou2023gptsgptsearlylook} can be found in \cref{tab:llm_exposure}} The resulting task assessments are then aggregated to construct three occupation level measures.

An occupation $j$ consists of $N_j$ tasks indexed by $i = 1, \ldots, N_j$. Each task is assigned a category $\omega_{ij} \in \{E0, E1, E2\}$ that denotes no exposure, direct exposure and potential exposure to \glsxtrshort{AI}, respectively. The first measure, denoted $\alpha_j$, captures directly exposed occupations and is defined as the directly exposed share of tasks to \glsxtrshort{AI} capabilities:

\[ \alpha_o = \frac{1}{N_o} \sum_{i=1}^{N_o} \mathbf{1}(\omega_{ij} = E1), \]

where $\mathbf{1}(\cdot)$ is a function equal to one if the condition holds and zero otherwise. The second measure, $\gamma_j$, captures the upper bound of exposure by weighting directly exposed and potentially exposed tasks equally:

\[ \gamma_j = \frac{1}{N_j} \sum_{i=1}^{N_j} \left[ \mathbf{1}(\omega_{ij} = E1) + \mathbf{1}(\omega_{ij} = E2) \right]. \]

This measure denotes the share of tasks in occupation $j$ that could be affected by \glsxtrshort{AI} capabilities without considering if any additional software is needed. Finally, $\beta_j$ provides an intermediate measure which reflect the uncertainty regarding potentially exposed tasks by assigning a half weight to them.

\[ \beta_j = \frac{1}{N_j} \sum_{i=1}^{N_j} \left[ \mathbf{1}(\omega_{ij} = E1) + 0.5\,\mathbf{1}(\omega_{ij} = E2) \right]. \]

All three measures are calculated using the task classifications produced by human raters and GPT-4o in \cite{martinez2024risk}. The result is six occupation level exposure measures between 0 and 1 that reflect the percentage of tasks in a given occupation which are exposed or potentially exposed to \glsxtrshort{AI} depending on the measure. $\alpha_j$ reflects direct exposure and is a lower bound, $\gamma_o$ represents an upper bound measure of potential exposure, and $\beta_o$ provides an intermediate benchmark. My estimations use all six measures allowing me to make results more robust.

\cref{fig:exposure_dist} plots the distribution across the three main exposure measures according to the minimum percent of exposed tasks in each occupation in the sample. It compares the results between the human assessed measures (\cref{fig:exposure_dist_human}) and the GPT assessed measures (\cref{fig:exposure_dist_gpt}). The first thing to notice is that across all measures by both assessments the distribution is skewed to the right, indicating low exposure in the majority of occupations. Meanwhile, the dashed lines denote the median exposure for each measure. In the GPT group, gamma exhibits the highest median exposure (5.9\%) and beta follows with 1.4\%. All other measures stay at 0\%.

Each treatment group is defined by the median exposure and the date ChatGPT released on the 4th quarter of 2022, specifically on  November 30th. Occupations above the threshold median classify as treated and occupations below serve as the control group. \cref{fig:occupations_and_workers_by_treat} plots the number of occupations in each treatment group in \cref{fig:occupation_by_treat}. The most restrictive measure is alpha ($\alpha$) because it classifies an occupation as exposed only by the ratio of directly exposed tasks. This measure is also more restrictive for the human measures (87) than the GPT measures (127). Another thing to note is that gamma and beta treatment group basically the same number of occupations even though they represent different exposure thresholds. This is a result of the majority of occupations having exposure 0. On the other hand, \cref{fig:workers_by_treat} highlights the economic importance of the treated groups by counting how much workers each group represents. Under the gamma measure, treated occupations account for 17.2 to 21.4 million workers. The alpha measure groups from 6.8 to 17.4 million workers. While this last count is smaller than the gamma count, it is in no way negligible.

\cref{fig:matrix_overlap} is the matrix that denotes the similarity between treatment groups by calculating in each cell the percentage of occupations present in the treatment group defined by the exposure measure on the x axis that are also included in the treatment group assigned by the measure in the y axis. There is an almost perfect overlap between the gamma and beta measures in both the human and the GPT assessment measures. Because of this, I keep only the beta and alpha exposure measures from this point onwards. Another thing to note is that the treatment group by the alpha measure is a subgroup of the beta measure, confirming the restrictiveness of the measure while making it directly comparable to the treatment beta group. The overall high percentage of overlap across groups also assures me that the human and GPT measures are grouping occupations appropriately with identifiable small differences.

Because one of my main outcomes is wages that are imputed, I want to make sure the imputation procedure is not driving my main results inaccurately. \cref{tab:imp_share_human} and \cref{tab:imp_share_gpt} report the share of imputed wages by quarter and treatment. The share of imputation is modest throughout the whole sample, ranging from 34.7\% to 48.9\% for the treated group and 23\% to 35.7\% for the control group in the alpha GPT measure. For the beta GPT measure the range for the treated is 34.2\% - 48.6\%, the control's range is 20.7\% - 32.9\%. The pattern is clearly increasing in the share of imputed wages and it is generally higher for the treated group. This result poses no threat to validity due to the result being in line with \cite{imput} as more people do not report their wage in the survey. Moreover, the share of imputed wages is increasing at a similar rate for the treatment and control groups. \cref{fig:imp_share_gap} plots the difference in the ratio of imputed wages between the treatment and the control groups across the main exposure measures. The fiigure displays a flat trend over time across the $\alpha$ and $\beta$ measures in both the human and GPT assessments. Treatment occupations having a higher share of imputed wages is also a common result due to the overall higher wage structure of exposed occupations.


\cref{tab:sumstats_exposure_alpha_gpt,tab:sumstats_exposure_beta_gpt,tab:sumstats_exposure_alpha_human,tab:sumstats_exposure_beta_human} display summary statistics for control and treated occupations by each exposure measure. Mean are weighted by occupation size. The before columns display the average outcome for the pre treatment period (2020Q3-2022Q4) and the after columns display the average outcome for the post treatment period (2023Q1-2026Q2). Across al four measures, exposed occupations differ from non exposed occupations over several core characteristics. Exposed groups are more skilled, have a higher ratio of female workers, have higher wages, are less informal, and are overall smaller occupations. Becuase of the higher formality, treated occupations are also more likely to be unionized and have access to social security. Although the overall differences are very similar across all four measures, the magnitude vaires slightly between them. The gap in schooling seems to be smaller in the treatment defined by human $\beta$ (12.2 versus 9.5 years) and GPT-4o Beta (12.8 versus 9.2) than under Human Alpha. Monthly earnings display a pattern off increases across both groups but it is not clear which is the net winner. Because the patterns shown by the summary statistics tell a very different basket of characteristics, comparisons in pure levels are not very useful. This motivates my approach to a \glsxtrfull{DID} estimation that will allow me to compare the evolution of outcomes between treated and control occupations.

Next, \cref{fig:wage_dist} plots the density of monthly earnings, observed and imputed, across occupations for the treated and control groups. All four distributions show what  \cref{tab:sumstats_exposure_alpha_gpt,tab:sumstats_exposure_beta_gpt,tab:sumstats_exposure_alpha_human,tab:sumstats_exposure_beta_human} showed before, that treated occupations are to the right on the non exposed occupations meaning the treated group has higher monthly earnings. The alpha measure shows the most distance between both distributions. This is consistent with the exposed group of akpha being a smaller subset of beta treatment group that has mostly directly exposed tasks.

Further motivitaing the design choice for the methodology, \cref{fig:trend_wage_imp,fig:trend_workers_imp,fig:trend_informal_imp} plot the evolution of the main outcomes for the treated and control groups across the sample. \cref{fig:trend_wage_imp} show a divergence in the treated and control groups after the treatment date. The divergence is more noticeable in panel (a) and (c) as these are the $\alpha$ measures showing direct exposure only. When plotting an index on the number of workers and informal workers in \cref{fig:trend_workers_imp,fig:trend_informal_imp}, there is no discerable effect that is homogeneous in all measures. Further evidence is required in order to coonclude an effect and that is why I turn to a \glsxtrshort{DID} estimation in the next section.





